\section{Discussion}

To summarize, our study demonstrates that achieving high thermodynamic fidelity in machine learning force fields is fundamentally bottlenecked by feature representations and training data composition, rather than raw model capacity. By evaluating models through active \gls{md} simulations and computing the \gls{kld} on Ramachandran marginals, we establish that static force prediction errors (MSE) alone are an insufficient criterion for model selection: models with comparable or even lower MSE can produce markedly different conformational distributions.

This work moves beyond a mere comparative evaluation of existing architectures by explicitly investigating the bounds of transferability within the highly combinatorial domain of dipeptides. Because amino acids can combine to form a vast chemical space, surrogate models must generalize to novel molecular graphs rather than simply interpolating within seen trajectories. By establishing a standard generalized Born (GB) implicit solvent model as a rigorous baseline, we revealed a critical limitation in representation strategies: while continuous partial charges \citep{machine_learning_implicit_solvation} successfully approximate the Boltzmann distribution for in-distribution data (KLD of 0.0939 versus the GB baseline's 0.1725), they fail to generalize, yielding a KLD of 0.1822 on unseen dipeptides---worse than the physics-based GB baseline (0.1289). This discrepancy exposes a fundamental deficiency in the transferability of continuous partial charge representations when confronted with unseen topologies.

In contrast, the topological and sequence-aware \textit{FF atom type + R} embedding \citep{klein_transferable_2024} achieves a KLD of 0.0165 on the out-of-distribution test dipeptide AA, surpassing the GB baseline by nearly an order of magnitude and outperforming the partial charge model by more than tenfold. Remarkably, this out-of-distribution performance (0.0165) even exceeds the model's in-distribution result on the validation dipeptide CI (0.0565), demonstrating that the learned representation does not merely interpolate within its training distribution but captures transferable physicochemical principles that generalize robustly to unseen molecular topologies. This constitutes, to the best of our knowledge, the first demonstration that a learned ML force field representation can systematically outperform a physics-based GB implicit solvent model on entirely unseen dipeptide systems, thereby establishing genuine zero-shot transferability within the combinatorial dipeptide space.

Furthermore, our ablation studies offer a nuanced understanding of the relationship between static force accuracy and dynamic sampling. During the feature embedding ablation, we observed a strong positive correlation ($r \ge 0.96$) between MSE and KLD, demonstrating that inadequate representations (such as baseline atom types) cause catastrophic force errors that inherently prevent accurate Boltzmann sampling. However, our trajectory length ablation reveals that once a model achieves a sufficient threshold of force accuracy with an optimal embedding, this correlation breaks down ($r = -0.41$ for AA). This confirms that while catastrophic static errors guarantee sampling failure, fine-grained differences in static MSE effectively decouple from the dynamic thermodynamic fidelity of the model. Consequently, explicit MD simulation and thermodynamic evaluation are indispensable for reliable model selection.

We also highlight the surprisingly strong sampling performance obtained from highly truncated training trajectories. When training on just the initial 5 ns of simulation data downsampled to 10{,}000 frames, the model achieves a KLD of 0.0232 on the test dipeptide AA---well below the GB baseline (0.1289) and competitive with the best-performing dataset configuration (0.0208 at 250 ns). Given that generating a 5 ns dataset is two orders of magnitude computationally cheaper than simulating the full 500 ns trajectory, this finding suggests that short-timescale structural data contains sufficient thermodynamic information to train accurate and transferable surrogate models, provided the feature embeddings are robust.

While our current study establishes a strong foundation using dipeptide systems, the combinatorial complexity of proteins necessitates further exploration. A primary avenue for future work will be evaluating these models on an already generated tetrapeptide dataset to explicitly test the transferability of our learned representations across varying sequence lengths. Additionally, given our core finding that feature embeddings---not parameter counts---are the primary bottleneck for thermodynamic fidelity, developing and testing novel embedding strategies remains a critical priority. Finally, combining these optimized, sequence-aware embeddings with more expressive, fully equivariant architectures, such as MACE \citep{batatia_mace_2022}, could further push the boundaries of accuracy, data efficiency, and zero-shot transferability in machine learning force fields.